\section{Conclusion}
\label{sec:conclusion}

We set out to understand why execution feedback enables code repair models to achieve 2.2$\times$ higher bug fix rates than global rewrites. Through controlled experiments isolating feedback signals from method confounds, we established a three-step causal chain: execution traces contain dense counterfactual information (84.7\% achieve CF-score $\geq$ 0.4), this information enables targeted edits (mean 4.2 lines versus 18.7 for binary feedback), and targeted edits achieve dramatically higher fix rates (68.4\% versus 31.2\%).

Our findings reframe the execution-versus-AI-feedback debate. The relevant distinction is not whether feedback comes from execution or an LLM, but whether feedback provides \emph{localization}. Detailed execution feedback achieves 100\% refinement success on failing problems; binary pass/fail achieves 60\%---same ground truth, different information structure. This suggests that well-designed AI critics providing accurate localization could, in principle, approach execution feedback performance.

Contrary to intuition, the execution advantage does not increase with task complexity. This task difficulty ceiling---where complex problems require more than localization to solve---defines a scope boundary for localization-based approaches and motivates future work on complementary strategies.

Our controlled comparison methodology---varying only feedback while holding method constant---provides a template for isolating active ingredients in refinement systems. We hope this approach enables principled optimization of code generation pipelines, moving beyond empirical trial-and-error to mechanistically grounded design.

Code and data are available at \url{https://github.com/anonymous/counterfactual-feedback}.
